% SGA 3，第 VI_B 篇，第一轮英文重建稿的中文翻译。
% 范围：引理10.9、段落10.10、命题10.11及其完整证明，含编者注95--97。
% 依据：Polo--Gille Exp6B-13oct24.pdf，本组件局部页71--73
% （合卷阅读版页412--414；再版页眉401--403；原版页边印刷页390--392）。
% 推论10.11.1始于局部页73，不在本组件内。
% OCR及Reinhold英文Markdown仅用作对照见证。
% 本限界组件不含实际图。

\medskip
\phantomsection\label{lem:sga3-VIB-10.9}%
\noindent\textbf{引理 10.9.---}
\emph{设 \(G_j\) 是有限呈示的 \(S_j\)-群，并令
\[
  G=G_j\times_{S_j}S.
\]
则 \(G^0\) 可表，当且仅当存在 \(i\geq j\)，使
\[
  (G_i)^0=(G_j\times_{S_j}S_i)^0
\]
可表。}

此条件是充分的，因为由 \hyperref[prop:sga3-VIB-3.3]{3.3}，函子
\(G\longmapsto G^0\) 与基变换交换。

反之，设 \(G^0\) 可表。则由 \hyperref[prop:sga3-VIB-3.9]{3.9}，
\(G^0\) 在 \(G\) 中开，且在 \(S\) 上拟紧，从而因 \(G\)
在 \(S\) 上有限呈示，它也在 \(S\) 上有限呈示。
由 \hyperref[lem:sga3-VIB-10.3]{10.3} 和
\hyperref[par:sga3-VIB-10.1]{10.1}，存在 \(i\geq j\) 及
\(G_i\) 的开群子概形 \(H_i\)，使
\[
  H_i\times_{S_i}S=G^0.
\]
结构态射 \(G^0\to S\) 的纤维几何连通
\textup{（\hyperref[prop:sga3-VIA-2.1.1]{VI\(_A\)，2.1.1}）}。
因此，由 EGA~IV\(_3\)，9.3.3和9.7.7，增大 \(i\) 后，可以假定
结构态射 \(H_i\to S_i\) 的纤维几何连通。%
\begingroup
\renewcommand{\thefootnote}{\fnsymbol{footnote}}%
\footnote[1]{\label{note:sga3-VIB-10.9-Hi-base}
译者注：法文源文印作 \(H_i\to S\)。由于 \(H_i\) 是 \(S_i\)-群，
且论证是将 \(G^0\to S\) 的连通性扩展到有限阶段，这里的底应为 \(S_i\)。}%
\endgroup\ 
由 \hyperref[lem:sga3-VIB-3.10.1]{3.10.1} 可知，\(H_i\) 的底空间
恰为 \((G_i)^0\)；因而 \(H_i\) 表示 \((G_i)^0\)。

\medskip
\phantomsection\label{par:sga3-VIB-10.10}%
\noindent\textbf{10.10.---}
回顾 \hyperref[par:sga3-VIB-10.0]{10.0} 所述情形的两个十分有用的特例
\textup{（参见 EGA~IV\(_3\)，8.1.2\textup{(a)}及\textup{(c)}）}：

\smallskip
\phantomsection\label{par:sga3-VIB-10.10-a}%
\noindent\textup{a)}
给定概形 \(X\) 的一点 \(x\)，令 \(S_0=\operatorname{Spec}\mathbb Z\)，
并考虑 \(x\) 的仿射开邻域所成的递减滤过逆向系 \((S_i)_{i\in I}\)。则
\[
  S=\operatorname{Spec}\mathcal O_{X,x}.
\]
特别地，若 \(x\) 是整概形 \(X\) 的泛点，则
\[
  S=\operatorname{Spec}\kappa(x).
\]

\smallskip
\phantomsection\label{par:sga3-VIB-10.10-b}%
\noindent\textup{b)}
令 \(S_0=\operatorname{Spec}\mathbb Z\)，并考虑仿射概形 \(S\) 的环的
有限生成 \(\mathbb Z\)-子代数所成的族 \((\mathcal A_i)_{i\in I}\)，
以包含关系赋予其预序。由于 \(\mathcal A_i\) 是 Noetherian 环，
在许多情况下，这使得我们可以由 Noetherian 情形过渡到一般情形。

现在给出关于 \hyperref[par:sga3-VIB-10.10-a]{(a)} 中所考虑的特例的两个结果。%
\footnote{\label{note:sga3-VIB-95}编者注：注意证明中也使用了情形\textup{(b)}。}

\medskip
\phantomsection\label{prop:sga3-VIB-10.11}%
\noindent\textbf{命题 10.11.---}%
\footnote{\label{note:sga3-VIB-96}编者注：对陈述作了简化，群子概形的情形
在\hyperref[cor:sga3-VIB-10.11.1]{推论10.11.1}中单独处理。}
\begingroup\itshape
设 \(S\) 是泛点为 \(\eta\) 的整概形。设 \(G\) 是 \(S\)-群，
\(Y,Z\) 是 \(S\)-概形，且它们都有限呈示。设
\[
  u,v,i:Y\longrightarrow G,
  \qquad
  j:Z\longrightarrow G
\]
是 \(S\)-概形态射。假设
\[
  i_\eta:Y_\eta\longrightarrow G_\eta,
  \qquad
  j_\eta:Z_\eta\longrightarrow G_\eta
\]
是闭浸入。

则存在 \(S\) 的非空开子集 \(S'\)，使经基变换得到的态射
\[
  i':Y'\longrightarrow G',
  \qquad
  j':Z'\longrightarrow G',
\]
是闭浸入，并使函子
\phantomsection\label{form:sga3-VIB-10.11-scheme-functors}%
\[
\begin{gathered}
  \underline{\operatorname{Transp}}(u',v'),\qquad
  \underline{\operatorname{Transp}}_{G'}
    \bigl(i'(Y'),j'(Z')\bigr),\\
  \underline{\operatorname{Transp}}^{\mathrm{str}}_{G'}
    \bigl(i'(Y'),j'(Z')\bigr)
\end{gathered}
\]
分别地，函子
\phantomsection\label{form:sga3-VIB-10.11-group-functors}%
\[
  \underline{\operatorname{Centr}}(u'),
  \qquad
  \underline{\operatorname{Norm}}_{G'}i'(Y'),
\]
由 \(G'\) 中在 \(S'\) 上有限呈示的闭 \(S'\)-子概形
\textup{（分别地，闭 \(S'\)-子群）}表示。
\par\endgroup

\emph{证明。}
我们将先在 \hyperref[par:sga3-VIB-10.10-a]{10.10\textup{(a)}} 的情形，
再在 \hyperref[par:sga3-VIB-10.10-b]{10.10\textup{(b)}} 的情形，
应用 \hyperref[par:sga3-VIB-10.1]{10.1} 的结果。
概形 \(G_\eta,Y_\eta,Z_\eta\) 在域 \(\kappa(\eta)\) 上平坦；
群 \(G_\eta\) 在 \(\kappa(\eta)\) 上分离
\textup{（\hyperref[sga3:VIA:0.3]{VI\(_A\)，0.3}）}；
而 \(i_\eta,j_\eta\) 是闭浸入。因此，由
\hyperref[par:sga3-VIB-10.1]{10.1}，存在 \(S\) 的仿射开子集
\(S'=\operatorname{Spec}(A')\)、\(A'\) 的 Noetherian 子环 \(A\)、
在 \(A\) 上平坦且有限呈示的 \(A\)-概形 \(G_A,Y_A,Z_A\)，以及态射
\[
  u_A,v_A,i_A:Y_A\longrightarrow G_A,
  \qquad
  j_A:Z_A\longrightarrow G_A,
\]
使得 \(G_A\) 是在 \(A\) 上分离的 \(A\)-群，\(i_A\) 和 \(j_A\)
是闭浸入，且
\[
  G\times_SS'=G_A\otimes_AA',
\]
其余概形和态射也有类似等式。
由于所考虑的 \(S'\) 上的函子由 \(\operatorname{Spec}(A)\) 上的
相应函子经基变换得到，只须对后者证明所述结果。

由 EGA~IV\(_2\)，6.9.2，将 \(A\) 替换为某个局部化 \(A_f\)
\textup{（从而将 \(S'\) 替换为仿射开子集 \(S'_f\)）}后，可以假定
\(G_A,Y_A,Z_A\) 在 \(A\) 上本质自由，其含义如
\hyperref[def:sga3-VIB-6.2.1]{6.2.1} 所述。%
\footnote{\label{note:sga3-VIB-97}编者注：事实上，由于 \(G_A\)
在 \(A\) 上平坦且有限呈示，它可被仿射开子集 \(G_1,\ldots,G_n\)
覆盖，使每个 \(\mathcal O(G_i)\) 都是平坦且有限呈示的 \(A\)-代数。
由 EGA~IV\(_2\)，6.9.2，存在 \(f_i\in A\)，使
\(\mathcal O(G_i)_{f_i}\) 是自由 \(A_{f_i}\)-模。
因此可以将 \(\operatorname{Spec}(A)\) 替换为仿射开子集 \(D(f)\)，
其中 \(f=f_1\cdots f_n\)，并对 \(Y_A\) 和 \(Z_A\) 同样处理。}
于是由 \hyperref[ex:sga3-VIB-6.2.4-b]{6.2.4\textup{(b)}} 和
\hyperref[ex:sga3-VIB-6.2.4-e]{6.2.4\textup{(e)}} 可知，在本命题的假设下，
所考虑的函子由 \(G_A\) 的闭 \(A\)-子概形表示。
由于 \(A\) 是 Noetherian 的，且 \(G_A\) 在 \(A\) 上有限呈示，
这些子概形在 \(A\) 上有限呈示；在
\(\underline{\operatorname{Centr}}(u_A)\) 和
\(\underline{\operatorname{Norm}}_{G_A}i_A(Y_A)\) 的情形中，
它们是 \(G_A\) 的闭 \(A\)-子群。
